\section{\amor{} jako instrument badawczy}

\subsection{Status narzędzia i własność intelektualna}

\amor{} jest autorskim rozwiązaniem autora niniejszej pracy, rozwijanym w ramach \bbg{}, firmy należącej do autora. To oznacza, że status narzędzia wymaga jawnego ujawnienia: model nie jest anonimowym zewnętrznym symulatorem, lecz infrastrukturą badawczą zaprojektowaną, zaimplementowaną i utrzymywaną przez autora jako projekt open-source. Z tego powodu szczególnie ważne są replikowalność, jawny plik zmian scenariusza BDP, dane Monte Carlo oraz możliwość krytyki założeń przez zewnętrznego czytelnika.

Taka pozycja narzędzia ma znaczenie metodologiczne, ale nie powinna być traktowana jako argument z autorytetu. Przeciwnie: skoro autor pracy jest zarazem autorem modelu, ciężar dowodu musi przejść na jawność mechanizmu, powtarzalność uruchomień i wyraźne ograniczenia interpretacyjne. Pytanie badawcze brzmi więc: co dzieje się z gospodarką polską, jeżeli do istniejącego modelu SFC-ABM zostanie dodany kanał BDP, a następnie zostanie on uruchomiony jako kontrfaktyczny eksperyment względem ścieżki bez BDP?

\amor{} operacjonalizuje intuicję ekonomii złożoności: zamiast zakładać reprezentatywnego agenta i racjonalne oczekiwania, konstruuje gospodarkę od dołu, z heterogenicznych agentów podejmujących ograniczenie racjonalne decyzje \parencite{arthur2021}. Hipoteza BDP nie jest więc w tej pracy wyłącznie argumentem narracyjnym. Zostaje zapisana w kodzie, uruchomiona jako kontrfaktyczny eksperyment i oceniona przez porównanie ścieżek modelu.

\subsection{Architektura i własności modelu}

\amor{} jest wykonywalnym modelem SFC-ABM gospodarki polskiej. Model obejmuje 10 000 heterogenicznych firm w sześciu sektorach, 100 000 agentów reprezentujących gospodarstwa domowe, sześć archetypów banków komercyjnych oraz agregat pozostałej części sektora bankowego. Obejmuje także bank centralny, sektor publiczny oraz sektor zewnętrzny z dynamicznym kursem PLN/EUR. Decyzje są podejmowane miesięcznie.

Punktem wyjścia eksperymentu jest produkcyjny punkt bazowy \amor{}, skalibrowany do stanu gospodarki Polski na dzień 2026-04-30. Praca nie kalibruje modelu od nowa, lecz wykorzystuje ten punkt startowy do przeprowadzenia kontrfaktycznych scenariuszy BDP.\footnote{Zakres i status tej powierzchni bazowej są udokumentowane w \path{docs/calibration-register.md}, \path{docs/scenario-registry.md}, \path{docs/empirical-validation-report.md} oraz \path{docs/empirical-validation/} repozytorium \amor{}.}

Liczba agentów nie jest próbą dosłownego odtworzenia 38-milionowej populacji Polski. Model wykorzystuje agentów jako skalowaną reprezentację heterogenicznych decyzji ekonomicznych, a agregaty makroekonomiczne są odnoszone do skali polskiej gospodarki przez mechanizmy kalibracyjne. Model należy więc interpretować jako narzędzie do badania kanałów transmisji i trajektorii makro-finansowych, a nie jako pełną mikrosymulację demograficzną każdego obywatela.

Z punktu widzenia tej pracy istotne są trzy własności modelu. Po pierwsze, przepływy są księgowo domknięte: transfer BDP nie może pojawić się wyłącznie jako dodatkowy dochód gospodarstw domowych, lecz musi mieć równoległy zapis po stronie wydatków publicznych, deficytu, długu lub finansowania. Po drugie, heterogeniczność firm, gospodarstw i banków pozwala obserwować nie tylko agregaty makro, ale także ograniczenia finansowania inwestycji i stabilność sektora bankowego. Po trzecie, scenariusze są replikowalne: każdy przebieg ma jawne ziarno losowania, horyzont miesięczny i artefakty wynikowe.

Szczegóły techniczne nie są częścią głównego argumentu ekonomicznego. Pełne macierze SFC, opis warstwy ledgerowej, granice formalnej weryfikacji i statusy kalibracji znajdują się w aneksach~\ref{app:sfc-matrices} i~\ref{app:engineering}. W tekście głównym ważny jest wniosek metodologiczny: każdy impuls polityki publicznej ma drugą stronę księgową, a odpowiedź firm zależy od bilansów, płynności i kosztu kapitału, nie tylko od presji na płace.

W tym sensie \amor{} pełni funkcję instrumentu badawczego, a nie ilustracji marketingowej. Jego wartość polega na tym, że hipotezę można zaimplementować, uruchomić, powtórzyć i sfalsyfikować, natomiast szczegółowy audyt techniczny zostaje oddzielony od głównej narracji.
